\documentclass[12pt,letterpaper]{article}

\usepackage[T1]{fontenc}
\usepackage[utf8]{inputenc}
\usepackage[spanish,mexico]{babel}
\usepackage{mathptmx} % Times New Roman
\usepackage[letterpaper,top=2.2cm,bottom=2cm,left=2.5cm,right=2.5cm]{geometry}
\usepackage{setspace}
\usepackage{parskip}
\usepackage{graphicx}

\setstretch{1.08}
\pagestyle{empty}

\begin{document}

\noindent Ciudad de México, a 2 de junio de 2026.

\vspace{0.6em}

\noindent \textbf{Comité de Selección}\\
\noindent \textbf{Programa de Veraneo Investigación y Tecnología 2026 (3.\textsuperscript{a} Edición)}\\
\noindent \textbf{INFOTEC}\\
\noindent \textbf{Presente.}

\vspace{0.8em}

Por este medio expreso mi interés en participar en el Programa de Veraneo Investigación y Tecnología 2026. Soy ingeniero en ciencias computacionales y maestrante en pedagogía; mi trabajo profesional y académico converge justo en el cruce que este programa propone: las tecnologías de la información analizadas desde lo técnico, lo ético y lo social.

Actualmente desarrollo una tesis de maestría sobre el análisis comparado de políticas públicas de educación en inteligencia artificial, con un estudio a profundidad de la \textbf{gobernanza de la inteligencia artificial en China} a lo largo de su trayectoria 2017--2025. Para ello construí un \textit{pipeline} de procesamiento de lenguaje natural que representa los documentos de política como \textit{embeddings} multilingües y mide su cercanía conceptual mediante búsquedas de similitud en espacios métricos, complementado con modelado de tópicos no supervisado. El programa me ofrece el acompañamiento experto que necesito para robustecer dos puntos críticos de esa metodología: el control lingüístico (distinguir señales reales de contenido frente a artefactos de traducción) y la validación de los hallazgos no supervisados frente a una codificación supervisada.

Aporto experiencia concreta y comprobable: diseño cursos de inteligencia artificial para Coursera en DeepLearning.AI, evalúo proyectos de procesamiento de lenguaje natural en la Maestría en Inteligencia Artificial Aplicada del Tecnológico de Monterrey y acumulo más de seis años enseñando ciencia de datos y programación. Estoy habituado a trabajar con rigor metodológico, a documentar de forma reproducible y a comunicar resultados técnicos a públicos diversos.

Mi compromiso es dedicar las 100 horas del programa a producir un protocolo de investigación sólido, replicable y publicable, que sirva tanto a mi tesis como a la comunidad interesada en el análisis cuantitativo de la política educativa digital y la gobernanza de la inteligencia artificial. Agradezco la oportunidad de ser considerado y quedo atento al proceso de entrevista.

\vspace{1.0em}

\noindent Atentamente,

\vspace{0.2em}

\noindent\includegraphics[height=2cm]{firma.png}\\[-0.35em]
\noindent\rule{6cm}{0.4pt}\\
\noindent \textbf{Hesus García Cobos}\\
\noindent hesusgc@gmail.com

\end{document}
